\chapter{Source coverage and audit index}
\label{chap:coverage-audit}

This chapter is a navigation and audit record, not an additional mathematical
input.  It maps the source-of-record snapshot to the chapters above and makes
the known source repairs visible.  Mathematical claims still live in their
individual nodes, with their own direct dependencies and status markers.

\section{Section-by-section coverage}

The introduction at \texttt{aimpaper/main.tex:78--252} is represented by the
classical Adams interfaces and low-line inputs in
Chapter~\ref{chap:classical-adams}, the permanent-cycle reduction in
Chapter~\ref{chap:near126}, and the conditional geometric endpoints in
Chapter~\ref{chap:final-kervaire}.  The two open questions remain open
proposition nodes and do not feed the proof graph.

Section 2, lines 253--669, is represented in
Chapter~\ref{chap:extension-ss}: the filtered two-term construction,
essential and inessential extensions, crossings, all comparison corollaries,
and the exact-at-the-middle theorem are separate nodes.  The regression
example is retained only where it fixes the intended crossing semantics.

Section 3, lines 670--997, is represented in the stable and synthetic
background chapters: the $H\mathbb F_2$-synthetic context, $\nu$, $\lambda$,
quotients, rigidity, $E_\infty$ formulas, synthetic lifts, and normalized
triangles all have explicit external or project status.

Section 4, lines 998--1272, is represented in the synthetic-extension
chapter.  It includes the $\lambda$-, $\rho$-, and $\delta$-extension spectral
sequences, the finite and infinite target formulas, quotient compatibility,
and the classical-crossing comparison.

Section 5, lines 1273--1601, is represented in the classical-page extension
chapter.  The $\widehat f_{r-1}$ page reduction, exact indeterminacy, crossing
comparison, $\lambda$-inversion, and the load-bearing Hopf-map regressions are
all explicit.

Section 6, lines 1602--2101, is represented in the generalized-rules chapter:
the generalized Leibniz rule, May's two-triangle input, the generalized
Mahowald trick, the page-loss theorem, and stretching to later pages are
separate obligations with proof routes.

Section 7, lines 2102--2782, is represented in
Chapter~\ref{chap:near126}.  Its graph separates BJM/BX literature, typed
high-stem evidence, choice-independence, the unique possible $d_{12}$,
Toda/two-extension arguments, the Hopf-cofiber obstruction, and the final
contradiction.  The source's circular-looking use at line 2496 is repaired by
first proving the filtration-leading-term Lemma
\ref{lem:toda-h0-leading-nonvanishing}; only after the Toda candidate
enumeration does Corollary~\ref{cor:h02x1259-d5-zero} deduce the unresolved
$d_5$ is zero.

The appendix at lines 2783--3290 and Table 1 in Section 7 are represented by
the twelve typed datasets in Chapter~\ref{chap:appendix-data}: 401 nonempty
rows and nine zero bands.  The active part of \texttt{aimpaper/112.tex}
defines Figure 9 only; Definition~\ref{def:near126-figure-contract} treats it
as a derived visualization, while its commented duplicate block is excluded.

\section{Recorded source repairs}

The Blueprint deliberately does not copy the following source defects:
\begin{enumerate}
  \item line 146's Hopf-survival inequality is read as $j\le3$, consistently
    with the following nonzero $d_2$ formula for $j\ge4$;
  \item line 604's class $z$ lies in $E_\infty(Z)$, not $E_\infty(Y)$;
  \item line 819's Bockstein variable is $\lambda$, not $\tau$;
  \item Proposition 4.6 separates the infinite target from its image in a
    finite $\lambda$-quotient;
  \item line 1522 uses $\widehat f_{r-1}$, consistent with Definition 5.1;
  \item Proposition 6.20 receives typed ranges for its correction classes;
  \item line 2496's backwards dependence is replaced by a noncircular
    filtration-leading-term argument;
  \item the lift from modulo $\lambda^3$ to modulo $\lambda^5$ in Lemma 7.20
    requires the independent obstruction-vanishing evidence node.
\end{enumerate}

\section{Status boundary}

Only exact declarations already compiled in the repository carry
\texttt{leanok}; only declaration names checked against the pinned Mathlib
revision carry \texttt{mathlibok}.  Every paper theorem, literature adapter,
data transcription, and missing background interface remains
\texttt{notready}.  Thus this document describes a complete formalization
plan and its proof frontier, not a completed Lean proof of the paper.
